\documentclass[10pt]{article}
\usepackage[T1]{fontenc}
\usepackage{lmodern,amsmath,amssymb,amsthm,microtype}
\usepackage[margin=0.92in]{geometry}
\usepackage[colorlinks=true,urlcolor=blue,citecolor=blue,linkcolor=blue]{hyperref}
\newtheorem{theorem}{Theorem}
\newtheorem{lemma}{Lemma}
\newtheorem{corollary}{Corollary}
\newcommand{\ord}{\operatorname{ord}}
\newcommand{\Cham}{\mathcal{C}}
\newcommand{\len}{\mathcal{L}}
\title{Every Prime Except 2 and 5 Divides All Rotations of Some Champernowne Prefix}
\author{Qichao Wang\\\small Hebei University of Technology, Tianjin, China}
\date{October 16, 2026}
\begin{document}
\maketitle
\begin{abstract}
Let $\Cham(m)$ denote the decimal integer obtained by concatenating
$1,2,\ldots,m$, and let $\len(m)$ be its number of decimal digits.
We prove that, for every prime $p\notin\{2,5\}$, some prefix $\Cham(m)$
and every cyclic digit rotation of that prefix are divisible by $p$.
This settles the positivity conjecture for OEIS A181373.
Our proof is elementary and constructive: for $p\ge7$ it supplies an
$m\le 10^{p-1}+p(p-1)-1$. Moreover, infinitely many such $m$ exist
for each admissible prime. A second construction in terms of the
multiplicative order of $10$ gives a sometimes sharper bound.
\end{abstract}

\section{Introduction}
Write
\[
  \Cham(m)=\underbrace{123456789101112\cdots m}_{\text{concatenated decimal blocks}},
  \qquad
  \len(m)=\sum_{k=1}^{m}\bigl(1+\lfloor\log_{10}k\rfloor\bigr).
\]
OEIS A007908 records $\Cham(m)$ and A058183 records $\len(m)$
\cite{oeis7908,oeis58183}. In A181373, the $n$-th term is the least
$m\ge1$ such that the $n$-th prime divides every cyclic \emph{digit}
rotation of $\Cham(m)$, or $0$ if no such $m$ exists
\cite{oeis181373}. Chai Wah Wu proposed in 2023 that the only zero
terms are those corresponding to $2$ and $5$.

The circular concatenations and their modular divisibility patterns
have been investigated by Rip\`a and Dalmasso \cite{ripa-dalmasso}
and later revisited by Rip\`a \cite{ripa-note}. We give a direct
existence proof by forcing the digit count into a prescribed congruence
class and then sweeping every residue of a prime field inside a single
fixed-digit-length block.

\section{A rotation criterion and an affine shift}
For a decimal integer $N$ of length $D$, cut its first $h$ digits into
an integer $A$, leaving a suffix $B$ of length $D-h$. Thus
$N=A10^{D-h}+B$. Moving the first $h$ digits to the end gives
$R_h=B10^h+A$, whence
\begin{equation}\label{eq:rotation}
  R_h=10^hN-A(10^D-1).
\end{equation}
This formula remains valid if the rotated digit string begins with zero:
its integer value is still $R_h$.
Consequently,
\begin{equation}\label{eq:criterion}
   p\mid N\quad\text{and}\quad p\mid(10^D-1)
   \quad\Longrightarrow\quad
   p\mid R_h\quad\text{for every }h.
\end{equation}
OEIS also records the equivalence between the full-rotation condition
and simultaneous divisibility in \eqref{eq:criterion}
\cite{oeis181373}; we only need sufficiency.

\begin{lemma}[Constant-step shift]\label{lem:shift}
Let $p\ge7$ be prime and let $r$ be a positive integer such that
$p\nmid r$ and $10^r\equiv1\pmod p$. Suppose $m,m+1,\ldots,m+r$
are all $d$-digit positive integers, where $d\equiv1\pmod r$.
Then
\begin{equation}\label{eq:shift}
  \Cham(m+r)-\Cham(m)\equiv-\frac{r}{9}\pmod p.
\end{equation}
\end{lemma}
\begin{proof}
Throughout the fixed-digit-length interval, concatenation satisfies
\[
  \Cham(t+1)=10^d\Cham(t)+(t+1).
\]
Since $d\equiv1\pmod r$ and $10^r\equiv1\pmod p$,
we have $10^d\equiv10\pmod p$. Iterating $r$ steps gives
\[
  \Cham(m+r)\equiv10^r\Cham(m)
           +\sum_{i=1}^{r}(m+i)10^{r-i}\pmod p.
\]
Set $A=\sum_{i=1}^r i10^{r-i}$.
As $10\not\equiv1\pmod p$ and $10^r\equiv1\pmod p$,
\[
  \sum_{i=1}^{r}10^{r-i}=\frac{10^r-1}{9}\equiv0\pmod p.
\]
A telescoping calculation yields
\[
  9A=\sum_{i=1}^{r}10^i-r
      =10\frac{10^r-1}{9}-r\equiv-r\pmod p.
\]
Here $9$ is invertible modulo $p$. Removing the term proportional
to $m$ proves \eqref{eq:shift}.
\end{proof}

\section{Resolution of the A181373 conjecture}
\begin{theorem}\label{thm:main}
For every prime $p\notin\{2,5\}$, there exist infinitely many
positive integers $m$ for which $p$ divides $\Cham(m)$ and all of
its cyclic digit rotations. No such $m$ exists for $p=2$ or $p=5$.
More quantitatively, for every prime $p\ge7$, such an $m$ can be chosen
with
\begin{equation}\label{eq:bound}
  m\le10^{p-1}+p(p-1)-1<10^p.
\end{equation}
Consequently, $\mathrm{A181373}(n)>0$ exactly when $n\notin\{1,3\}$.
\end{theorem}
\begin{proof}
Fix a prime $p\ge7$, and put
\[
   r=p-1,\qquad d=p,\qquad L=10^{p-1},\qquad
   H=\len(L-1).
\]
By Fermat's little theorem, $10^r\equiv1\pmod p$;
furthermore $d\equiv1\pmod r$ and $p\nmid r$.
Choose the unique $\delta\in\{0,1,\ldots,r-1\}$ satisfying
\begin{equation}\label{eq:delta}
   \delta\equiv-H-1\pmod r,
\end{equation}
and consider $p$ endpoints
\begin{equation}\label{eq:candidates}
   m_j=L+\delta+jr,\qquad j=0,1,\ldots,p-1.
\end{equation}
All $m_j$ have exactly $d=p$ decimal digits: indeed,
\[
   0\le\delta+jr\le(r-1)+(p-1)r=p(p-1)-1
                         <9\cdot10^{p-1}.
\]
(The last inequality is immediate for $p\ge7$.)
Consequently, by counting all $d$-digit blocks between $L$ and $m_j$,
\[
  \len(m_j)=H+d(m_j-L+1)
            =H+p(\delta+jr+1)\equiv H+\delta+1\equiv0\pmod r.
\]
It follows that
\begin{equation}\label{eq:length}
   10^{\len(m_j)}\equiv1\pmod p.
\end{equation}

Applying Lemma~\ref{lem:shift} successively inside this interval,
using $r=p-1$, shows
\[
   \Cham(m_j)\equiv\Cham(m_0)-\frac{jr}{9}
                  \equiv\Cham(m_0)+\frac{j}{9}\pmod p.
\]
As $j$ varies from $0$ to $p-1$, the right side takes every residue
modulo $p$ exactly once. In fact, choose the unique
\begin{equation}\label{eq:jchoice}
   j\equiv-9\Cham(m_0)\pmod p,\qquad 0\le j<p.
\end{equation}
For this $j$, we obtain $p\mid\Cham(m_j)$, while
\eqref{eq:length} and \eqref{eq:criterion} imply that $p$ divides
all digit rotations of $\Cham(m_j)$. The bound \eqref{eq:bound}
follows from \eqref{eq:candidates}.

To obtain infinitely many distinct witnesses, use instead
$d=1+t(p-1)$ for each integer $t\ge1$, with $L=10^{d-1}$.
The construction above applies without change, and the $d$-digit
interval accommodates the same $p$ candidate endpoints.
As $d$ strictly increases, the witnesses are distinct.

For $p=3$, take $m=2$: $\Cham(2)=12$, $\len(2)=2$,
and $3\mid12$ as well as $3\mid(10^2-1)$.
Indeed, every $m\equiv0$ or $2\pmod3$ yields
$\Cham(m)\equiv m(m+1)/2\equiv0\pmod3$ because $10\equiv1\pmod3$;
thus there are infinitely many witnesses for $3$ as well.

Finally, for $p=2$ or $5$, the first decimal digit of $\Cham(m)$
is $1$. A digit rotation that moves this digit to the last place
ends in $1$ and hence cannot be divisible by $2$ or $5$.
This completes the proof.
\end{proof}

\begin{corollary}[An order-sensitive bound]\label{cor:order}
Let $p\ge7$ be prime, and write $r=\ord_p(10)$. There exists a
witness $m$ with
\[
  m\le10^{2r}+pr-1<10^{2r+1}.
\]
Thus one may take the smaller of this bound and \eqref{eq:bound}.
\end{corollary}
\begin{proof}
Set $d=2r+1$, $L=10^{2r}$, $H=\len(L-1)$, and choose
$\delta\in\{0,\ldots,r-1\}$ with $\delta\equiv-H-1\pmod r$.
Put $m_0=L+\delta$ and consider $m_j=m_0+jr$ for $0\le j<p$.
Every candidate has digit count divisible by $r$, since
\[
  \len(m_j)=H+d(\delta+jr+1)\equiv H+\delta+1\equiv0\pmod r.
\]
Moreover, $\delta+(p-1)r\le pr-1$; because $p\mid 10^r-1$,
we have $p<10^r$, hence $pr<r10^r<9\cdot10^{2r}$.
Thus all candidates remain $d$-digit integers.

The step in Lemma~\ref{lem:shift} is now $-r/9$, not
$1/9$ as in the special case $r=p-1$. In particular,
\[
  \Cham(m_j)\equiv\Cham(m_0)-jr\,9^{-1}\pmod p.
\]
Since $r=\ord_p(10)\mid p-1$, the inverse $r^{-1}\pmod p$
exists. Choose explicitly
\begin{equation}\label{eq:order-choice}
  \boxed{j\equiv 9\,r^{-1}\Cham(m_0)\pmod p,\qquad 0\le j<p.}
\end{equation}
Then $p\mid\Cham(m_j)$, while $r\mid\len(m_j)$ implies
$p\mid10^{\len(m_j)}-1$. The rotation criterion
\eqref{eq:criterion} completes the construction.
\end{proof}

\section{Discussion}
The argument does not determine the least witness $m$ in A181373;
it proves its existence, gives an explicit (generally nonminimal)
construction, and establishes infinitely many valid endpoints for
each admissible prime. The only external tools are elementary
modular arithmetic and Fermat's little theorem.

\begin{thebibliography}{9}
\bibitem{oeis181373} OEIS Foundation Inc.,
\emph{A181373: Least $m$ such that a prime divides a prefix and its cyclic rotations},
\url{https://oeis.org/A181373}, accessed October 16, 2026.
\bibitem{oeis7908} OEIS Foundation Inc.,
\emph{A007908: Concatenation of the positive integers},
\url{https://oeis.org/A007908}, accessed October 16, 2026.
\bibitem{oeis58183} OEIS Foundation Inc.,
\emph{A058183: Number of digits in the concatenation of the first $n$ positive integers},
\url{https://oeis.org/A058183}, accessed October 16, 2026.
\bibitem{ripa-dalmasso} M. Rip\`a and E. Dalmasso,
Patterns related to the Smarandache circular sequence primality problem,
\emph{Notes on Number Theory and Discrete Mathematics},
\textbf{18} (2012), no. 1, 29--48.
\url{https://nntdm.net/volume-18-2012/number-1/29-48/}.
\bibitem{ripa-note} M. Rip\`a,
A note on my 2012 paper on patterns related to the Smarandache circular sequence primality problem,
Zenodo (2025), \url{https://doi.org/10.5281/zenodo.17382562}.
\end{thebibliography}
\end{document}
